\section{Prescribing linear information before testing circuit size}
Suppose $s\sim cN/n$, with fixed $0<c<1$. A rank-$r$ matrix $A$ specifies $r$ parity checks on a table. Its syndrome is the vector $Ax$ of answers. For a fixed syndrome $y$ we may treat $Ax=y$ as known in advance and ask what remains of MCSP on that affine space. Neither the matrix nor the choice of coordinates on its solution space need be easy to compute.

\begin{theorem}
Fix $0<\rho<c<1$ and $r=\lfloor\rho N\rfloor$. Uniformly for every rank-$r$ matrix $A$, each fiber $Ax=y$ contains an affine space of easy tables of dimension at least $(c-\rho-o(1))N$. One may use the same direction space for all $y$ for this $A$. If the syndrome is uniform, then with probability $1-o(1)$ the following statements hold:
\begin{enumerate}
\item The largest affine space of easy tables in the fiber has dimension $(c-\rho+o(1))N$, and the largest affine space of hard tables there has dimension $(1-c+o(1))N$.
\item The number of easy tables in the fiber has logarithm $(c-\rho+o(1))N$. The bounded-error parity query complexity of MCSP restricted to this fiber, and the logarithm of its bounded-error minimum parity tree size measured in leaves, are both $(c-\rho+o(1))N$.
\end{enumerate}
If instead $c<\rho\le1$, the number of attainable syndromes of easy tables has logarithm $(c+o(1))N$, for every $A$. A uniform syndrome is attainable with probability $2^{-(\rho-c+o(1))N}$.
\end{theorem}
\begin{proof}
For any fixed $0<\epsilon<c-\rho$, extend the $r$ rows to a linearly independent list of $r'=\lfloor(c-\epsilon)N\rfloor$ rows, denoted by $A'$. The block construction of Section~2 yields one linear image $U$ of dimension $r'$ that projects surjectively under $A'$, with every word computable in $(r'+o(N))/n\le s$ gates. The common direction $U\cap\ker A$ has dimension $r'-r$. Intersection of $U$ with each fiber is a coset of this space. Sending $\epsilon$ slowly to zero gives the first statement.

Let $M_y=|A_s\cap\{Ax=y\}|$. By the circuit census,
$\mathbb E_y M_y=|A_s|2^{-r}\le 2^{(c-\rho+o(1))N}$.
Markov's inequality and the preceding lower bound then give $\log M_y=(c-\rho+o(1))N$ with high probability. The same upper bound applies to the dimension of any easy flat in that fiber. Within such a fiber, any rank-$k$ internal parity check splits it into $2^k$ cosets. If $2^k>M_y$, one of these cosets consists entirely of hard tables. Take $k=\lceil\log(M_y+1)\rceil$. Its dimension is $N-r-k=(1-c-o(1))N$. The global hard-flat upper bound in the preceding section applies to all such cosets.

For the query lower bound use the common random linear encoder of the companion manuscript \emph{Small Boolean truth tables under adaptive parity tests}, Theorem~2, at rank $r'$. The ensemble is independent of all queried rows, its failure on any fixed rank-$r'$ system is $o(1)$, and every message maps to an easy table. Given $A,y$, draw an encoder conditioned on full projection onto $A$, then draw a uniform message among those mapping into the fiber. For any fixed additional $q\le r'-r$ tests independent modulo $A$, extend the joint list to rank $r'$. With conditional probability $1-o(1)$ the encoder has full projection on the list, since the conditioning event has probability $1-o(1)$. For each such encoder all $2^q$ additional syndromes are equiprobable at the prescribed $y$. Each affine cell of codimension $q$ inside the fiber thus has at least $(1-o(1))2^{-q}$ mass under this law. The error bound is uniform over cells. A parity tree with at most $2^{r'-r-\epsilon' N}$ leaves assigns at most $2^{-\epsilon'N}$ of the fiber's uniform mass to leaves of codimension above $r'-r$. Summing the smaller masses over the other leaves shows that the two transcript laws have total variation $o(1)$, even after averaging over independent tree coins. In a typical fiber the uniform hard probability tends to one by the bound on $M_y$. Correct bounded-error recognition therefore requires larger trees. A depth-$q$ tree has at most $2^q$ leaves, and $\lceil\log M_y\rceil+2$ random parity checks suffice by the fingerprint argument. Taking the fixed margins to zero proves the two query conclusions in probability.

Finally, when $\rho>c$, the circuit census bounds the number of syndromes by $2^{(c+o(1))N}$. Restrict $A$ to any $\lfloor(c-\epsilon)N\rfloor$ of its rows. Every partial syndrome is attained by an easy table by Section~2, giving at least $2^{(c-\epsilon)N-O(1)}$ distinct complete syndromes. This proves the last assertion.
\end{proof}

The typical-syndrome qualification matters. For instance an affine space of dimension less than $(c-\epsilon)N$ supported on an initial segment of addresses can consist entirely of easy tables, by the lookup bound of Section~2. If $c+\rho<1$, however, the fraction of uniform tables that are easy tends to zero in \emph{every} fiber by the global count. In this regime the parity query lower bound applies in every fiber.
